% UNINTEGRATED WHOLE-B18 WORKING APPENDIX, 24 October 2026.
% Same existing general compact paper family; not a standalone paper.
% Source-only transcription in progress. Remaining bodies are listed at end.
% No compilation, worker import, test or scientific execution has occurred.

\section{Terminal walls in a whole \(B_{18}\) diagram}
\label{app:wholeB18}
Fix the original strict interior unit point \(z\) and an original
farthest vertex \(p\), and assume that the entire incident diagram
is \(B_{18}\). Normalize the actual old normals by
\[
 n_i=(e_i-e_{i+1})/\sqrt2\quad(1\leq i\leq17),\qquad n_{18}=e_{18}.
\]
Write
\[
 p^2=-1,\qquad z=E(p-x),\qquad E>0,\quad x\perp p,\quad
 x_1>\cdots>x_{18}>0,\quad \norm x^2<1.
\]
For \(P_i=e_1+\cdots+e_i\), put
\[
 S_i=\ip x{P_i},\qquad f_i=\sqrt i-.58/\sqrt2,\qquad
 d_K=\sqrt K-\sqrt{K-1}.
\]
The actual duals are \(\omega_i=\sqrt2P_i\) for \(i\leq17\)
and \(\omega_{18}=P_{18}\). In particular
\(\ip{\omega_i}{\omega_j}=2\min(i,j)\) for \(i,j\leq17\).
An index is strong when \(S_i>f_i\); equality belongs to the weak case.
There are no old spectator directions. All subsequent references
in the original frame concern this one vector \(x\).

We use the actual positive projection and future-point pressure lemmas
already proved in the main text. Their hypotheses will be checked for
each point below. A positive correction on actual old normals leaves
every unknown wall inequality nonpositive by acuteness. A spacelike
corrected normal is used only inside its positive finite projection;
it is never normalized as a chamber point.

\subsection{The early window and actual head pressures}
\label{app:B18bulk}
Deleting \(i\), \(5\leq i\leq14\), retains
\(A_{i-1}+B_{18-i}\), with both ranks at least four.
The finite one-node list consists of isolation; one left end bond
of label three or four; a left \(D\) or permitted \(E\) position;
one right long-end simple bond; or an end-to-end bridge.
Two contacts on one retained component create a cycle. A joined
rank-eighteen diagram is classical; exceptional types add no case.
The right \(B_3\) at \(i=15\) additionally allows its short-end
simple bond, and the right \(B_2\) at \(i=16\) allows either end.
These last roles will be checked separately.

For an actual edge with \(C=\cosh L=qR\), \(q=\sqrt{2i}\),
clearing the positive denominators of its directional pressure
shows that
\[
 R>T(q):=\frac{50}{29}
    -\frac{q-29/50}{q(2q-29/50)}
 \quad\Longrightarrow\quad S_i>f_i.
\]
The function increases for \(q\geq\sqrt{10}\), and substitution gives
\[
 T(q)<83/50\ (q<7),\qquad
 T(q)<801/500\ (q<15/4),\qquad
 T(q)<159/100\ (q<16/5).
\]
The finite-white Schur calculations have the following complete
non-helper bounds; \(\phi=(1+\sqrt5)/2\).
\[
\begin{array}{c|l}
\text{role}&\text{lower bound for }R\\ \hline
\text{left near three, own at least five}
 &[i(\phi+1)-1]/\sqrt{2i(i+1)}\\
D\text{ position}&(i\phi/2+1)/\sqrt{2i}\\
E\text{ position}&2\phi/2+3/8\\
\text{right long three, positive remaining own}
 &1+1/\sqrt2\\
\text{right }B_3\text{ short three, positive own}
 &1+1/\sqrt2\\
\text{end-to-end bridge}&(i+1)/\sqrt{2i}.
\end{array}
\]
The near-three bound increases with \(i\). At \(i=6\) it is
\((6\phi+5)/\sqrt{84}>801/500\), which covers \(6,7\);
at \(i=8\) it exceeds \(83/50\), covering larger \(i\).
At \(i=5\), every own at least seven gives \(R>5/3\).
For the \(D\) row, \(i=5\) gives \(R>159/100\), and
\(i\geq6\) gives \(R\geq2927/1750>83/50\).
An \(E\) position has \(i\leq8\) and retained residual at most
\(1/4\). The bridge residual is \(1/(2i)\); its numerator is
at least \((i+1)/(2i)\), proving the displayed lower bound,
including own zero. Left label-four roles have the affine
small-own alternative or the larger direct edge bounds:
near own three has \(C=(i-1)+i/\sqrt2\).
The positive-null and impossible-time boundaries in the common
prefix inventory apply to the actual retained blocks at these
ranks. For the right long-three own-zero boundary one can also
see the kernel directly: use coefficient one on the cap and
old \(i\), two on old \(i+1,\ldots,17\), and \(\sqrt2\) on
old short \(18\). Every selected product vanishes.

An isolated finite cap has the actual finite/core pressures used
in the earlier component lemmas. At own five its loss is less than
\(25/76+5/24<.58\); own seven has loss
\(1000/4779+9/25<.57\); the remaining finite interval has its
largest loss at \(i=5\), own eight:
\(\sqrt{10}+\sqrt2-4<.578\).
For an isolated dotted cap put \(u=g^{-2}\). Its positive
reciprocal core gives
\[
 \ip x{\omega_i}>\sqrt{2i-u}-\sqrt{u(1-u)}.
\]
For \(i\geq6\) the loss is at most
\[
 a_i u+\sqrt{u(1-u)},\qquad
 a_i=\frac1{\sqrt{2i}+\sqrt{2i-1}}<37/250,
\]
whose maximum is less than \(.5795\), since
\(\sqrt{1+(37/250)^2}<1011/1000\).
Thus only an isolated dotted cap at \(5\) may remain weak.

There is one early near-three exception, own five at \(5\).
For \(g=\phi/2\), the actual cap has displacement
\(v=g\omega_5+\omega_4/2\). On the right actual \(H_4\), correct by
\[
 w=g(3n_5/2+n_6+n_7/2),\qquad
 \delta=(5-3\phi)/8>0.
\]
The corrected normal is spacelike and its finite projection is
actual; the extra old \(4\) product is negative, and all unknown
wall signs follow by positivity. Its pure spatial displacement is
\[
 Y=(1/2+3g/4)\omega_4+(g/4)\omega_8,\qquad
 A^2=(7+13\phi)/2,\quad Y^2=(33+49\phi)/8.
\]
The finite-point pressure gives
\[
 \ip xY>18/5,\qquad Y^2<281/20,\qquad
 \sum Y_j<43/5.
\]
Independently, the earlier rank-valid \(.61\) prefix lemma supplies
a strong \(K=11\) or \(12\). Its old floors exceed \(72/25\)
or \(303/100\). The vector \(Y+(3/4)P_K\) therefore has
squared demand exceeding squared norm by \(.0401\) or
\(.78625625\), respectively, a contradiction.
This use of the earlier lemma concerns \(B_{18}\geq B_{16}\)
and precedes any new \(.58\) or strong-five conclusion.

Consequently every weak ordinary site is a pure far helper
\[
 c_i=A_i p-g\sqrt2P_i-e_1/\sqrt2,\qquad
 A_i^2=2ig^2+2g-\tfrac12,\qquad
 \phi/2\leq g<6/5,
\]
apart from the isolated dotted \(5\).
The upper cutoff follows from
\(R^2\geq4g-2\geq14/5>(83/50)^2\) at \(g\geq6/5\).

For adjacent pure helpers the exact time/dual calculation gives
\[
 \ip{c_i}{c_{i+1}}
 \geq \frac12-g_i g_{i+1}
           +\frac14(g_i/g_{i+1}+g_{i+1}/g_i)
 \geq1-g_i g_{i+1}>-11/25.
\]
Their actual native gap forces orthogonality. The positive two-cap
replacement, with every other old wall retained, has positive
inverse Schur matrix; its weighted loss is less than
\(7/16+1/(4\sqrt{10})<.58\).
Thus both sites cannot be weak. This uses their actual new
parameter range, rather than inclusion in an older weak class.
In particular some \(K=14\) or \(15\) is strong.

The same actual helper feet give \(x_1>.5455\) at a weak
helper \(5\). For an isolated dotted weak \(5\), its reciprocal
core gives
\[
 U'=P_5+(u/2)(e_6-e_5),\qquad
 \ip x{U'}>T'>1.8185,\qquad .4<u<1.
\]
Indeed the loss is bounded by
\(a u+\sqrt{u(1-u)}\), \(a=\sqrt{10}-3<13/80\);
for \(u\leq.4\) it is less than \(.565<.58\), and globally
it is less than \(.588\).
The pressure bound follows from
\(\sqrt{10}>3.162,\sqrt2<99/70\).
These are actual point pressures, not assigned floors at weak nodes.

\subsection{Raw references, head isolation and the uniform paid prefix}
\label{app:B18references}
A reference \(w\) has support through a strong \(K\), sum \(f_K\),
actual pressure \(\ip xw>A_w\), and effective energy
\(L=2A_w-w^2\). Point insertion need not leave \(L=w^2\).
For a new actual pressure \(\ip xV\geq G\), weight \(\mu\geq0\),
and \(Y=V-P_K\), the completed-square rule is
\[
 \norm x^2>L+2\mu(G-f_K-\ip wY)-\mu^2Y^2,
 \qquad \mu\leq\beta_K .
\]
The inequality is strict when the old strong reference retains
a positive coefficient. Every use below checks its individual prices.

Let \(d_i=\sqrt{i+1}-\sqrt i\).
The first-five energy exceeds \(.6668157\), and rational lower
numerators for \(10^7d_i^2\), \(i=5,\ldots,16\), are
\[
 455488,385186,333704,294372,263340,238230,
 217494,200080,185248,172465,161332,151549.
\]
For a direct verification, put \(a=4i+2\) and use
\(d_i^2(a-d_i^2)=1\), whence
\(d_i^2>(a^2-1)/(a(a^2-2))\).
Thus the unomitted energies satisfy
\[
 E_{14}>.9241299,\quad E_{15}>.9413764,\quad
 E_{16}>.9575096,\quad E_{17}>.9726645.
\]
A missing knot \(j\geq6\) loses less than
\(1/[32(j-1)^3]\); nonadjacency makes these losses disjoint.
Their total is less than
\((1/32)\int_4^\infty t^{-3}\,dt=1/1024\).
Each terminal capacity is at least \(d_K\).

Before any head payment, the three initial-five cases give the
following full contracts:
\[
\begin{array}{c|cc|c}
\text{case at five}&L_{14}>&L_{15}>&w_1<\\ \hline
\text{strong}&.9231&.9403998&.3652\\
\text{pure helper}&.95255&.96895&.5455\\
\text{isolated dotted weak}&.9291&.9455&.377 .
\end{array}
\]
In the strong case interpolate through all actual strong floors,
skipping every weak knot. The first mean is \(f_5/5<.3652\).
For a helper, interpolate through \((1,.5455),(6,f_6),(K,f_K)\).
Its first-six energy exceeds
\(.5455^2+(f_6-.5455)^2/5>.74385\);
the tails to \(14,15\) exceed \(.2087,.2251\).
All slopes are positive decreasing.

For isolated dotted weak \(5\), choose strong \(J=6\) or \(7\)
and a reference constant first \(J\), then straight to \(K\).
Insert \(\mu[U'-(5/J)P_J]\), with \(\mu=.21,.2\).
The residual squares are bounded by
\[
 5-u+u^2/2-25/J<77/150,\ 194/175
\]
and the deficits exceed \(.119,.221\).
Its individual \(P_J\) prices are \(.175,1/7\),
below the capacities \(>.1783,>.1628\) already at \(K=14\).
The gains exceed \(.027342,.044057\).
The bare first-\(J\) energies exceed \(.69312,.7139\);
the tails to \(14\) exceed \(.2087,.1715\) and to \(15\)
exceed \(.2251,.18825\).
The full first coordinates are less than
\[
 .3399+.21(1-5/6)=.3749,\qquad
 .31938+.2(1-5/7)<.377.
\]
This uses the straight reference for these prices. A curved
reference with a smaller \(P_6\) capacity is not substituted.

The actual long-head opposite is isolated. Deleting old \(1\)
retains \(B_{17}\); its finite enlargement is isolation or a
simple bond to its long endpoint \(2\). In the latter case
\[
 v=g\omega_1+\omega_2/2,\qquad A^2=2g^2+2g,\qquad C_1=2g+1.
\]
Own zero has zero time; own three is an actual affine triangle.
Every other own gives \(x_1>\sqrt{\sqrt2-1}>.643\).
If \(5\) is strong, replace the first-five mean \(<.3652\)
by this pressure, gaining
\((5/4)(.643-.3652)^2>.0964\) over \(L_{14}>.9231\).
If \(5\) is weak, choose strong \(J=6,7\) and interpolate
through \(1,J,K\); the first energies exceed \(.8034,.8361\)
and the tails exceed \(.2087,.1715\).
Each case exceeds one, including isolated dotted weak \(5\).
Therefore a head has time \(A_H^2=2g_H^2-1\).
Owns zero, three, four have impossible time; six and parallel
have the positive actual \(6/3\) and parallel nulls.
Only own five, finite own at least seven, and dotted own remain.

For later reference, every non-five head has a point pressure
\(\ip xZ\geq h\), with \(\sum Z_j=1\) and the following
full local data:
\[
\begin{array}{c|c|c|c}
\text{own}&Z&h&R=Z^2\\ \hline
\text{finite }\geq7&
 (1-u/2)(e_1+e_2)+(u-1)e_3&
 [\sqrt{2-u}-\sqrt{7u-4-3u^2}]/\sqrt2&
 3-4u+3u^2/2\\
\text{dotted}&(1-u/2)e_1+(u/2)e_2&
 [\sqrt{2-u}-\sqrt{u(1-u)}]/\sqrt2&
 1-u+u^2/2 .
\end{array}
\]
Here \(u=g_H^{-2}\); the first row has
\(1<u\leq\cos(\pi/7)^{-2}\), the second \(0<u<1\).
Their actual positive cores are
\(c_H+g_H^{-1}n_1+2(g_H^{-1}-g_H)n_2\)
and \(c_H+g_H^{-1}n_1\), respectively.
Their cap product is zero, the stated old-root products are
nonpositive and all unknown signs follow termwise.
They are future, giving these pure spatial pressures.
Both obey
\begin{equation}\label{app:eq:B18headF}
 .8h-.16R>.2893.
\end{equation}
For native seven, \(.9009<g_H<.901\) gives \(h>.4315,R<.349\),
and the left side exceeds \(.28936\).
For finite labels at least eight, \(1<u<1.172\).
Writing \(D=(u-1)(4-3u)\), the derivative is
\[
 \frac{.8}{\sqrt2}\left(-\frac1{2\sqrt{2-u}}
                 -\frac{7-6u}{2\sqrt D}\right)+.64-.48u.
\]
It is negative on \((1,7/6]\); on \([7/6,1.172)\),
\(\sqrt D>.284\), its positive term is \(<.032\)
and its other terms total less than \(-.20\).
At own eight \(h>.4388,R<.373\), giving \(>.29136\).
For dotted \(u\), the two chords
\(\sqrt{2-u}\geq1.4142-.379u\) on \([0,1/2]\) and
\(\sqrt{2-u}\geq1.4494-.4494u\) on \([1/2,1]\)
give, respectively,
\[
 .8h-.16R>
 .63987152-.0944u-.5657\sqrt{u(1-u)}>.30587152,
\]
\[
 .8h-.16R>
 .69978064-.2142u-.5657\sqrt{u(1-u)}>.29023064.
\]
Use
\(bu+c\sqrt{u(1-u)}\leq(b+\sqrt{b^2+c^2})/2\),
with radical bounds \(.5736,.6049\).
This proves the entire finite and unbounded dotted head seal.

+
\subsection{One uniform paid prefix, without a terminal premise}
\label{app:B18paidprefix}
If the head has own five, its actual unit normal is
\[
 c_H=A_Hp-(\phi/\sqrt2)e_1,\qquad A_H^2=(\phi-1)/2.
\]
A weak pure helper at \(i=5,6,7\) would have mutual magnitude
\[
 M_i(g)=A_H\sqrt{2ig^2+2g-\tfrac12}
             -\frac\phi2(2g+1).
\]
This increases in \(g,i\) on the allowed range and
\(M_7(6/5)<0\), since the positive squared comparison has
difference \((814\phi-1392)/100<0\).
An isolated dotted weak \(5\) has own \(h<2\), by its edge
pressure and \(f_5^2<7/2\). Its increasing mutual magnitude
\[
 A_H\sqrt{10h^2-1}-\phi h
\]
lies strictly between zero and \(1/4\) for \(1<h<2\):
the endpoint squared comparisons are
\((7\phi-11)/2>0\) and \(377/16-(29/2)\phi>0\).
Both cases contradict actual acuteness or the native gap.
Thus a native-five head makes \(5,6,7\) strong independently
of any terminal hypothesis.

The head's actual finite \(H_4\) projection gives
\[
 \ip x{P_4/4}\geq h_0,\qquad
 h_0^2=(\phi-1)/4,\qquad h_0>.39307.
\]
At strong \(K=14,15\) the curved reference, omitting every
nonadjacent internal weak knot, has energies
\[
 w_{14}^2>.924016575,\qquad w_{15}^2>.94125107.
\]
Indeed the complete raw energies above lose at most
\((.0112^2+.008^2+.0061^2)/2=.000113325\)
for \(K=14\), and \(<.00012533\) for \(K=15\);
these bounds include every permitted hole.
The \(P_5\) capacity is \(>.151\).
Pay weight \(.56\) on \(P_4/4-P_5/5\), costing \(.112\).
Its square is \(1/20\), it is orthogonal to the raw reference,
and the single-square gain exceeds
\[
 1.12(.39307-.3652)-.56^2/20=.0155344.
\]
Hence this initial/head case gives
\[
 L_{14}>.939550975,\qquad L_{15}>.95678547.
\]

All other initial/head cases give larger uniform energies.
For a non-five head and strong \(5\), pay \(.4(Z-P_5/5)\).
Its price is \(.08<.151\), and
\eqref{app:eq:B18headF} gives gain
\(.2893-.8(.3652)+.032>.02914\).
For a weak helper \(5\), the raw contracts above already
give \(L_{14}>.95255,L_{15}>.96895\), with no further head graft.

For isolated dotted weak \(5\), choose strong \(J=6,7\) and
a curved reference to \(K\), constant first \(J\).
Its first mean, individual capacity, and the square of
\(Y_5=U'-(5/J)P_J\) obey
\[
\begin{array}{c|ccc}
J&m<&\beta_J>&Y_5^2<\\ \hline
6&.3399&.1435&77/150\\
7&.31938&.1366&194/175 .
\end{array}
\]
The raw first-\(J\) energies exceed \(.69317,.71400\);
the complete tails to \(14\) exceed \(.2117654,.1732468\)
after omission losses \(<.000267,.000181\).
Thus the raw energies exceed \(.9046684,.8870658\).
At \(15\), the additional increment is \(>.0172465\).
Let \(Y_H=Z-P_J/J\). The full cross is
\(\ip{Y_H}{Y_5}=1-5/J\), while
\(\ip w{Y_H}=\ip w{Y_5}=0\).
The one vector \(w+.4Y_H+.09Y_5\) spends \(.85/J\),
below each displayed capacity. Its entire gain is
\[
 (.8h-.16R)-.8m+.16/J
 +.18(T'-5m)-.0081Y_5^2-.072(1-5/J).
\]
It exceeds \(.04930,.06699\), respectively. The resulting
\(K=14\) energies exceed \(.9539684,.9540558\), and the
\(K=15\) energies exceed \(.9712149,.9713023\).
The earlier \(.21/.2\) insertion is not also charged.

In every case, therefore, there is one paid reference \(V\),
supported through a strong \(K=14\) or \(15\), with sum \(f_K\),
unspent terminal capacity \(\beta_K\geq d_K\), and
\begin{equation}\label{app:eq:B18uniformprefix}
 2A_V-V^2>L_K,\qquad
 L_{14}=.939550975,\quad L_{15}=.95678547 .
\end{equation}
All corrections have zero sum and support through seven.
Their individual nonnegative coefficients were checked before
the terminal pressure is used.

We shall also need the following independent implication:
if \(17\) is strong, the head has own five.
For a non-five head and strong \(5\), the curved reference to
\(17\) has energy \(>.97133\), and its head gain \(>.02914\)
already exceeds one. If \(5\) is a helper, the full curved
reference gains \(>.74385-.7126=.03125\) in its first six
coordinates, yielding
\(E_{17}+.03125-1/1024>1.0029\).
If \(5\) is isolated dotted, the analogous curved references
to \(17\) have energies \(>.9530,.9354\); the joint
head/isolated-five gains just proved exceed \(.04930,.06699\).
Again the totals exceed one. Thus every non-five head is
incompatible with strong \(17\). This argument precedes any
classification of its terminal opposite.

\subsection{Unconditional terminal sums and isolated opposites}
\label{app:B18sumcutoffs}
For \(N=16,17,18\), let \(Y=P_N-P_K\).
It has support after \(K\), square \(N-K\), and zero product
with the entire corrected reference in
\eqref{app:eq:B18uniformprefix}.
Paying the nonnegative weight \(d_K\leq\beta_K\) shows that
\(S_N\geq G\) would imply
\[
 \norm x^2>L_K+2d_K(G-f_K)-(N-K)d_K^2.
\]
The rational root bounds
\[
\begin{gathered}
3.60555<\sqrt{13}<3.60556,\quad
3.74165<\sqrt{14}<3.74166,\\
3.87298<\sqrt{15}<3.87299
\end{gathered}
\]
give
\(.1361<d_{14}<.1362,\ .1313<d_{15}<.1314\).
Every bound follows by positive squares and
\(d_K=1/(\sqrt K+\sqrt{K-1})\).

For \(S_{17}\geq3.760\), the extra terms are bounded below by
\[
 2(.1361)(.428)-3(.1362)^2=.06085028,\qquad
 2(.1313)(.297)-2(.1314)^2=.04346028.
\]
Their totals exceed \(1.000401255,1.00024575\).
For \(S_{18}\geq3.830\), the increments are
\(.06135384,.04457632\), giving totals
\(>1.000904815,>1.00136179\).
Thus
\begin{equation}\label{app:eq:B18sumbounds}
 S_{17}<3.760,\qquad S_{18}<3.830,\qquad
 C_{17}<10.879,\quad C_{18}<9.81 .
\end{equation}
The edge cosines use the actual directions
\(P_N/\sqrt N\):
\(C_N<(N+G_N^2)/(N-G_N^2)\).
The two fractions are
\(31.1376/2.8624<10.879\) and
\(32.6689/3.3311<9.81\).
These are simultaneous consequences for the same \(x\),
not energies added from different references.

An entirely isolated \(17\)-opposite has
\[
 c=A p-\sqrt2gP_{17},\qquad A^2=34g^2-1>0.
\]
Own zero has negative time. Own three, four, six and parallel
have positive actual null combinations, respectively,
\[
 c+n_{16}+2n_{17}+\sqrt2n_{18},\qquad
 c+\sqrt2n_{17}+n_{18},\qquad
 g c+n_{17}+n_{16}/2,\qquad c+n_{17}.
\]
Selected products vanish, cap time is positive, and all
unselected wall products are nonpositive. Compactness forbids them.
For own five, correct on the actual \(H_4\) by
\(g(3n_{17}/2+n_{16}+n_{15}/2)\), with positive residual
\(\delta=(5-3\phi)/8\). The actual finite projection gives
pressure on \(V=P_{14}+\tfrac34(e_{15}+\cdots+e_{18})\)
strictly greater than \((A-\sqrt\delta)/(\sqrt2g)\).
Ordering implies \(\ip xV\leq S_{17}\).
Since \(A^2=(17\phi+15)/2\), the rational bounds
\(A>4.61,\sqrt\delta<.136,\sqrt2g<1.145\) give
\(S_{17}>(4.61-.136)/1.145>31/8\).
Every finite own at least seven has the positive core
\(c+g^{-1}n_{17}+2(1-g^2)g^{-1}n_{16}\).
Its norm is \(-(1-g^2)(4g^2-3)/g^2\), selected cap/own
products vanish, and the other selected product is negative.
All unknown wall signs follow from positive actual coefficients.
With \(u=g^{-2}\in(1,4/3)\), divide its displacement by
\(\sqrt2g\). Its coordinates are one through fifteen,
\(2-u\) at sixteen, and \(u/2\) at seventeen and eighteen.
Its product with \(x\) is at most \(S_{17}\) by ordering.
The actual pressure is at least
\(\sqrt{17-u/2}-\sqrt{(7u-4-3u^2)/2}\).
The first radical exceeds \(4.04\) and the second is at most
\(1/\sqrt{24}<.2042\); hence \(S_{17}>34/9\).
All these finite cases contradict the same prefix rule:
the gains at \(G=34/9\) exceed \(.065,.047\).

For every dotted \(g>1\), put \(u=g^{-2}\).
The actual future reciprocal core \(c+n_{17}/g\)
has norm \(-(1-u)\). Its displacement is
\(U=\sqrt2gP_{17}-n_{17}/g\).
For \(r=17-K\), choose
\[
 Y=U-\sqrt2gP_K,\qquad
 \nu=d_K/(\sqrt2g),\qquad
 Y^2=2r/u-2+u.
\]
The residual is supported after \(K\), so is orthogonal
to the entire already paid reference. Its price is
\(\nu\sqrt2g=d_K\), allowing equality at the terminal.
The gain is exactly
\[
 G_K(u)=\sqrt2d_K\!
 \left[\sqrt{34-u}-\sqrt{2K}+.58-\sqrt{u(1-u)}\right]
 -d_K^2(r-u+u^2/2).
\]
The chord \(\sqrt{34-u}\geq\sqrt{34}-.0864u\) follows from
\(\sqrt{34}>5.8309,\sqrt{33}>5.7445\).
The circle square gives
\(\sqrt{u(1-u)}\leq1/4+u-u^2\).
Using \(\sqrt{28}<5.2916,\sqrt{30}<5.4773\),
the bracket exceeds \(u^2-1.0864u+.8693\) or
\(u^2-1.0864u+.6836\), both positive.
Also \(\sqrt2d_{14}>.19247,\sqrt2d_{15}>.18568\).
Use the upper \(d_K\) only in the negative constant and
the lower \(d_K\) in \(d_K^2(u-u^2/2)\).
This gives
\[
\begin{split}
G_{14}(u)&>.183208395u^2-.190576198u+.111662851,\\
G_{15}(u)&>.177060155u^2-.184483062u+.092398928.
\end{split}
\]
For \(F=au^2-bu+c\), the identity
\(4a(F-\tau)=(2au-b)^2+4a(c-\tau)-b^2\)
has positive last constants
\[
 \frac{808151227201376}{10^{18}}\quad(\tau=.061),\qquad
 \frac{244086609159516}{10^{18}}\quad(\tau=.044).
\]
Thus the full dotted interval gives gains \(>.061,>.044\),
and totals \(>1.000550975,>1.00078547\).
No isolated \(17\)-opposite survives.

% Further proof bodies remain: terminal inventory, finite/dotted far helper,
% near-six and distinct/common terminal neighbors. No integration yet.

